\section{Numerical solution of the quantitative model}\label{app:numerical_solution}

This appendix summarizes the quantitative model in Section~\ref{sec:quantitative} in a form that is directly usable for computation.
The objective is to solve for equilibrium policy and pricing functions in a heterogeneous-beliefs economy with Epstein--Zin preferences, production with hiring chosen one period in advance, and a stationary wealth distribution generated by Uzawa discounting.

\paragraph{Objects to compute.}
Given the state \(X_t\), the key endogenous objects are the risk-neutral expectation of productivity growth \(\mathcal{L}_t\), the surplus-claim price-dividend ratio \(n_t\), the consumption-claim price-dividend ratios \(\{n_{i,t}\}_{i=1}^I\), and consumption shares \(\{s_{i,t}\}_{i=1}^I\) (equivalently, Pareto weights \(\{\lambda_{i,t}\}\)).
All other quantities (SDFs, returns, hours, wages, and output) follow from these objects.

In contrast to the simpler baseline model in the main text, which is conveniently described in terms of wealth shares, the quantitative solution tracks heterogeneity using Pareto weights (or, equivalently, consumption shares). This representation is convenient for the perturbation-based solution.

\paragraph{Computation.}
An exact closed-form solution is unavailable away from log preferences.
We therefore compute the solution using a \textit{third-order perturbation} around the non-stochastic (deterministic) steady state, as described in Section~\ref{sec:quantitative}.
The third-order approximation is indispensable to generate time variation in expected returns and other risk-premium objects.

Concretely, we take a Taylor expansion of the equilibrium conditions below around the steady state and solve for the coefficients of the resulting policy rules for the endogenous variables as functions of the state.
We then simulate the implied nonlinear state-space system using the approximated decision rules to compute moments and impulse responses. For simulation, we use a \textit{pruned} state-space representation of the higher-order solution to avoid spurious explosive dynamics induced by truncation of the perturbation expansion (see \citealt{andreasen2018pruned}).

\subsection{Equilibrium conditions}

\paragraph{Aggregate productivity growth.} Aggregate log productivity growth, $\hat{x}_t \equiv \log x_t$, follows the following process under the objective measure:
\begin{equation}
    \hat{x}_t = \underbrace{\mu_x + \rho_x (\hat{x}_{t-1}-\mu_x)}_{\overline{x}_{t-1}} + \sigma_{x,t-1}\epsilon_{x,t}.
\end{equation}
Aggregate productivity follows the following process under the beliefs of investor $i$:
\begin{equation}
    \hat{x}_t =  \underbrace{\mu_{x,i} + \rho_{x,i} (\hat{x}_{t-1}-\mu_{x,i}) + v_{i,t-1}}_{\overline{x}_{i,t-1}} + \sigma_{x,i,t-1}\epsilon_{x,i,t},
\end{equation}
where the belief shock $v_{i,t}$ follows the process:
\begin{equation}
    v_{i,t} = \rho_{v,i} v_{i,t-1} + \sigma_{v,i} \epsilon_{v,i,t}.
\end{equation}
\noindent Following \citet{bansal2004risks}, we allow for stochastic volatility:
\begin{equation}
    \sigma_{x,t}^2 = q_t,
    \qquad
    \sigma_{x,i,t}^2 = q_t\frac{\overline{\sigma}_{x,i}^2}{\overline{\sigma}_x^2},
    \qquad i=1,\ldots,I,
\end{equation}
where \(q_t\) follows an AR(1) process in levels:
\begin{equation}
    q_t = \bar{q} + \rho_q(q_{t-1}-\bar{q}) + \sigma_q \epsilon_{q,t}, \qquad \epsilon_{q,t}\sim\mathcal{N}(0,1).
\end{equation}
In the baseline specification, we consider the special case with constant volatility: \(q_t\equiv \bar q\), in which \(\sigma_{x,t}=\overline{\sigma}_x\) and \(\sigma_{x,i,t}=\overline{\sigma}_{x,i}\).
In a robustness exercise, we allow for stochastic volatility, where the volatility shocks \(\epsilon_{q,t}\) may be correlated with productivity and belief shocks, with \(\text{Corr}(\epsilon_{x,t},\epsilon_{q,t})=\rho_{xq}\) and \(\text{Corr}(\epsilon_{v,t},\epsilon_{q,t})=\rho_{qv}\). 

\paragraph{Likelihood ratios.} We assume that subjective beliefs are absolutely continuous with respect to the true distribution, so for any random variable $g_{t+1}$ we can write:
\begin{equation}
    \mathbb{E}_{i,t}[g_{t+1}] = \mathbb{E}_{t}[\ell_{i,t+1} g_{t+1}],
\end{equation}
where $\ell_{i,t+1}$ denotes the Radon--Nikodym derivative of the subjective probability density with respect to the true probability. We also assume that all shocks are Gaussian, so we can write this derivative as the likelihood ratio of two Gaussians:
\begin{equation}
    \ell_{i,t+1} = \frac{\sigma_{x,t}}{\sigma_{x,i,t}}
    e^{-\frac{(\hat{x}_{t+1} - \overline{x}_{i,t})^2}{2\sigma_{x,i,t}^2} + \frac{(\hat{x}_{t+1} - \overline{x}_t)^2}{2\sigma_{x,t}^2}},
\end{equation}
where \(\sigma_{x,t}\) and \(\sigma_{x,i,t}\) denote the objective and subjective conditional standard deviations of \(\hat{x}_{t+1}\) given information at time \(t\).

\paragraph{Preferences: EZ+ Uzawa.} To ensure that a non-degenerate stationary wealth distribution exists in this economy, we assume that investors have external Uzawa preferences, such that the discount factor for household $i$ is given by:
\begin{equation}
    \beta_{i,t} = \beta e^{-\kappa s_{i,t}},
\end{equation}
where $s_{i,t} \equiv \log \frac{\tilde{C}_{i,t}}{Y_t A_t}$ denotes the relative consumption of type $i$, and $Y_t \equiv h_t^\alpha - \xi e^{-\hat{x}_t} \frac{h_t^{1+\nu}}{1+\nu}$ denotes detrended net output.
When $\kappa > 0$, the discount factor is decreasing in the consumption share of agent $i$, so households become more impatient as their consumption increases. 

The stochastic discount factor for household $i$ is given by
\begin{equation}
    \Lambda_{i,t+1} = \beta_{i,t}^\theta \left( \frac{\tilde{C}_{i,t+1}}{\tilde{C}_{i,t}} \right)^{-\frac{\theta}{\psi}} \left( \frac{n_{i,t+1} + 1}{n_{i,t}} \frac{\tilde{C}_{i,t+1}}{\tilde{C}_{i,t}} \right)^{\theta-1},
\end{equation}
using the fact that $R_{i,n,t+1} = \frac{N_{i,t+1}}{N_{i,t} - \tilde{C}_{i,t}} = \frac{n_{i,t+1}+1}{n_{i,t}} \frac{\tilde{C}_{i,t+1}}{\tilde{C}_{i,t}}$, where $n_{i,t} \equiv \frac{N_{i,t}}{\tilde{C}_{i,t}}-1$ represents the price-dividend ratio on a claim on household $i$'s consumption.

We can then write the SDF as follows:
\begin{equation}
    \Lambda_{i,t+1} = \beta_{i,t}^\theta \left( e^{\Delta s_{i,t+1} + \Delta y_{t+1} + \hat{x}_{t+1}} \right)^{-\gamma} \left( \frac{n_{i,t+1}+1}{n_{i,t}} \right)^{\theta-1},
\end{equation}
where $y_{t} \equiv \log Y_t$. 

\paragraph{Consumption claim's price-dividend ratio.} The pricing condition for household's portfolio return is given by
\begin{equation}
    1 = \mathbb{E}_{i,t}\left[ \Lambda_{i,t+1} R_{i,n,t+1} \right]. 
\end{equation}
We can write the expression above as follows:
\begin{equation}
    1 = \mathbb{E}_{t}\left[ \ell_{i,t+1} \beta^\theta e^{-\theta \kappa s_{i,t}} \left( e^{\Delta s_{i,t+1} + \Delta y_{t+1} + \hat{x}_{t+1}} \right)^{1-\gamma} \left( \frac{n_{i,t+1}+1}{n_{i,t}} \right)^{\theta} \right].
\end{equation}

\paragraph{Optimal risk sharing.} The optimal risk sharing condition can be written as follows:
\begin{equation}
    \ell_{i,t+1} \Lambda_{i,t+1} = \ell_{1,t+1} \Lambda_{1,t+1}.
\end{equation}

We can write this condition as follows:
\begin{equation}
    \ell_{i,t+1}\beta_{i,t}^\theta \left( \frac{\tilde{C}_{i,t+1}}{\tilde{C}_{i,t}} \right)^{-\gamma} \left( \frac{n_{i,t+1} + 1}{n_{i,t}} \right)^{\theta-1}=
    \ell_{1,t+1}\beta_{1,t}^\theta \left( \frac{\tilde{C}_{1,t+1}}{\tilde{C}_{1,t}} \right)^{-\gamma} \left( \frac{n_{1,t+1} + 1}{n_{1,t}} \right)^{\theta-1},
\end{equation}

It is useful to work with a recursive formulation of the problem. 
Let $\hat{\lambda}_{i,t}$ denote the Pareto weight of household $i$, and write the optimal risk sharing condition as follows:
\begin{equation}\label{eq:optimal_rs}
    \hat{\lambda}_{i,t-1} \ell_{i,t} \tilde{C}_{i,t}^{-\gamma} (n_{i,t}+1)^{\theta-1} = \hat{\lambda}_{1,t-1} \ell_{1,t} \tilde{C}_{1,t}^{-\gamma} (n_{1,t}+1)^{\theta-1}.
\end{equation}
where $\hat{\lambda}_{i,t}$ denote the \textit{Pareto weight} on household $i$, which evolves according to the law of motion:
\begin{equation}\label{eq:optimal_rs2}
    \hat{\lambda}_{i,t} = \zeta_t \hat{\lambda}_{i,t-1} \beta_{i,t}^\theta \ell_{i,t} \left( \frac{n_{i,t}+1}{n_{i,t}} \right)^{\theta-1},
\end{equation}
where $\zeta_t$ is a scaling process. 

In the special case of CRRA preferences, $\theta =1$, constant discount factor, $\kappa = 0$, and rational expectations, $\ell_{i,t} = 1$, relative Pareto weights are constant and the optimal risk sharing condition boils down to the expression with time-separable preferences: $\hat{\lambda}_i \tilde{C}_{i,t}^{-\gamma} = \hat{\lambda}_1 \hat{C}_{1,t}^{-\gamma}$.

Taking the ratio of Eq.~\eqref{eq:optimal_rs} at $t+1$ and $t$, we obtain
\begin{equation}
    \frac{\hat{\lambda}_{i,t}}{\hat{\lambda}_{i,t-1}} \frac{\ell_{i,t+1}}{\ell_{i,t}} \left( \frac{\tilde{C}_{i,t+1}}{\tilde{C}_{i,t}} \right)^{-\gamma} \left( \frac{n_{i,t+1}+1}{n_{i,t}+1} \right)^{\theta-1} =
    \frac{\hat{\lambda}_{1,t}}{\hat{\lambda}_{1,t-1}} \frac{\ell_{1,t+1}}{\ell_{1,t}}\left( \frac{\tilde{C}_{1,t+1}}{\tilde{C}_{1,t}} \right)^{-\gamma} \left( \frac{n_{1,t+1}+1}{n_{1,t}+1} \right)^{\theta-1},
\end{equation}
which coincides with $\ell_{i,t+1} \Lambda_{i,t+1} = \ell_{1,t+1} \Lambda_{1,t+1}$.
We can choose $\hat{\lambda}_{i,-1}$ such that condition~\eqref{eq:optimal_rs} is satisfied at $t=0$.

From Eq.~\eqref{eq:optimal_rs} and the market clearing for goods, $\sum_{j=1}^I \mu_j \tilde{C}_{j,t} = A_t Y_t$, we can solve for relative consumption:
\begin{equation}
    \frac{\hat{C}_{i,t}}{A_t Y_t} = \frac{\hat{\lambda}_{i,t-1}^{\frac{1}{\gamma}}\left[ \ell_{i,t} (n_{i,t}+1)^{\theta-1}\right]^{\frac{1}{\gamma}} }{\sum_{j=1}^I \mu_j \hat{\lambda}_{j,t-1}^{\frac{1}{\gamma}}\left[ \ell_{j,t} (n_{j,t}+1)^{\theta-1}\right]^{\frac{1}{\gamma}}}.
\end{equation}
Taking logs, we can write the expression above as follows:\small
\begin{equation}
    s_{i,t} - s_{1,t} = \lambda_{i,t-1} - \lambda_{1,t-1} + \frac{1}{\gamma} \left[ \log \frac{\ell_{i,t}}{\ell_{1,t}} + (\theta-1) \log \frac{n_{i,t}+1}{n_{1,t}+1} \right],
\end{equation}\normalsize
where $\lambda_{i,t} \equiv \frac{1}{\gamma}\log \hat{\lambda}_{i,t} $, and $\sum_{i=1}^I \mu_i e^{s_{i,t}} =1$. We can solve for $s_{1,t} $ in terms of $\tilde{s}_{i,t} \equiv s_{i,t} - s_{1,t}$, $i = 2,\ldots, I$, using the condition
\begin{equation}
    s_{1,t} = - \log\left( \mu_1 + \sum_{i=2}^I \mu_i e^{\tilde{s}_{i,t}} \right). 
\end{equation}

The law of motion of $\lambda_{i,t}$ is given by
\begin{equation}
    \lambda_{i,t} = \lambda_{i,t-1}  + \frac{1}{\gamma}\left[ - \theta \kappa s_{i,t}+\log \ell_{i,t} + (\theta-1) \log (1+n_{i,t}^{-1}) + \log \beta^\theta \zeta_t\right].
\end{equation}
We are free to choose the value for $\zeta_t$. For instance, we can set $\zeta_t = \beta^{-\theta}$, so the last term inside brackets above is equal to zero. Alternatively, we can set $\zeta_t$ such that $\lambda_{1,t} = 0$. 

\paragraph{Asset pricing.} The price of a risk-free bond is given by
\begin{equation}
    P_{f,t} = \mathbb{E}_t\left[ \ell_{i,t+1} \Lambda_{i,t+1} \right],
\end{equation}
and the risk-free rate is given by $R_{b,t} = P_{f,t}^{-1}$. 

Let $R_{r,t+1} $ denote the return on a claim on total surplus, or net output $Y_t$. Let $n_{t}$ denote the price-dividend ratio on the surplus claim, we then write the return as follows:
\begin{equation}
    R_{r,t+1} = \frac{n_{t+1}+1}{n_t} e^{\Delta y_{t+1} + \hat{x}_{t+1}}. 
\end{equation}
The pricing condition for the surplus claim is given by
\begin{equation}
    1 = \mathbb{E}_t\left[ \ell_{i,t+1} \Lambda_{i,t+1} R_{r,t+1} \right].
\end{equation}


\paragraph{Production.} The key endogenous object \(\mathcal{L}_t\) is the risk-neutral expectation of productivity growth. It evolves according to:
\begin{equation}
    \mathcal{L}_t = R_{b,t} \mathbb{E}_t\left[ \ell_{i,t+1} \Lambda_{i,t+1} e^{\hat{x}_{t+1}} \right]
    \qquad \Rightarrow \qquad
    \mathbb{E}_t\left[ \ell_{i,t+1} \Lambda_{i,t+1} (e^{\hat{x}_{t+1}}- \mathcal{L}_t) \right] = 0.
\end{equation}

Hours and wages are given by:
\begin{equation}
    h_{t} = \left( \frac{\alpha \mathcal{L}_{t-1}}{\xi} \right)^{\frac{1}{1+\nu-\alpha}}, \qquad\qquad
    w_t = \xi \left( \frac{\alpha \mathcal{L}_{t-1}}{\xi} \right)^{\frac{\nu}{1+\nu-\alpha}}. 
\end{equation}

Output is given by
\begin{equation}
    e^{y_t} = h_t^\alpha - \xi e^{-\hat{x}_t} \frac{h_t^{1+\nu}}{1+\nu}.
\end{equation}

\subsection{Summary of equilibrium conditions}

\paragraph{Roadmap.} The equilibrium system can be organized as:
\begin{itemize}
    \item \textbf{State dynamics:} the exogenous laws of motion for \(\hat{x}_t\), \(q_t\), and \(v_{i,t}\), plus endogenous evolution of \(\mathcal{L}_t\) and Pareto weights \(\lambda_{i,t}\).
    \item \textbf{Risk sharing and pricing:} the consumption-claim Euler equations for each type and the surplus-claim pricing equation.
    \item \textbf{Production block:} hours, wages, and detrended net output \(Y_t\).
\end{itemize}

\paragraph{State dynamics:}\small
\begin{align}
    \hat{x}_t &= \mu_x + \rho_x (\hat{x}_{t-1}-\mu_x) + \sigma_{x,t-1}\epsilon_{x,t}\\
    q_t &= \bar{q} + \rho_q(q_{t-1}-\bar{q}) + \sigma_q \epsilon_{q,t}\\
    v_{i,t} &= \rho_{v,i} v_{i,t-1} + \sigma_{v,i} \epsilon_{v,i,t}, \qquad i = 1,\ldots, I\\
        \mathcal{L}_t &= R_{b,t} \mathbb{E}_t\left[ \ell_{1,t+1} \Lambda_{1,t+1} e^{\hat{x}_{t+1}} \right]\\
            \lambda_{i,t} &= \lambda_{i,t-1} + \frac{1}{\gamma}\left[- \theta \kappa (s_{i,t}-s_{1,t}) + \log \frac{\ell_{i,t}}{\ell_{1,t}} + (\theta-1) \log \frac{1+n_{i,t}^{-1}}{1+n_{1,t}^{-1}} \right], \qquad i = 2,\ldots, I
\end{align}\normalsize
where $\lambda_{1,t} = 0$. State vector: $X_t = (\hat{x}_t, q_t, \mathcal{L}_t,\{v_{i,t}\}_{i=1}^I, \{\lambda_{i,t}\}_{i=2}^I) $.  

\paragraph{Equilibrium conditions:}\small
\begin{align}
        % \overline{x}_{i,t} &= \mu_{x,i} + \rho_{x,i} (\hat{x}_{t-1}-\mu_{x,i}) + v_{i,t-1}\\
        %     l_{i,t+1} &= \frac{\sigma_x}{\sigma_{x,i}} e^{-\frac{(\hat{x}_{t+1} - \overline{x}_{i,t})^2}{2\sigma_{x,i}^2} + \frac{(\hat{x}_{t+1} - \overline{x}_t)^2}{2\sigma_x^2}}\\
            1 &= \mathbb{E}_{t}\left[ \ell_{i,t+1} \beta^\theta e^{-\theta \kappa s_{i,t}} \left( e^{\Delta s_{i,t+1} + \Delta y_{t+1} + \hat{x}_{t+1}} \right)^{1-\gamma} \left( \frac{n_{i,t+1}+1}{n_{i,t}} \right)^{\theta} \right], \qquad i= 1, \ldots, I\\
            s_{i,t} - s_{1,t} &= \lambda_{i,t-1}+ \frac{1}{\gamma} \left[ \log \frac{\ell_{i,t}}{\ell_{1,t}} + (\theta-1) \log \frac{n_{i,t}+1}{n_{1,t}+1} \right], \qquad i = 2, \ldots I\\
            1 &= \sum_{i=1}^I\mu_i e^{s_{i,t}}\\
            R_{b,t}^{-1} &= \mathbb{E}_{t}\left[ \ell_{1,t+1} \Lambda_{1,t+1} \right]\\
            1 &= \mathbb{E}_t\left[ \ell_{1,t+1} \Lambda_{1,t+1} \frac{n_{t+1}+1}{n_t} e^{\Delta y_{t+1} + \hat{x}_{t+1}} \right]\\
            \overline{R}_{r,t} &= \mathbb{E}_t\left[ \frac{n_{t+1}+1}{n_t} e^{\Delta y_{t+1} + \hat{x}_{t+1}} \right]\\
            h_{t} &= \left( \frac{\alpha \mathcal{L}_{t-1}}{\xi} \right)^{\frac{1}{1+\nu-\alpha}}\\
             w_t &= \xi \left( \frac{\alpha \mathcal{L}_{t-1}}{\xi} \right)^{\frac{\nu}{1+\nu-\alpha}}\\
             e^{y_t} &= h_t^\alpha - \xi e^{-\hat{x}_t} \frac{h_t^{1+\nu}}{1+\nu},
\end{align}\normalsize
where\small
\begin{equation}
    \Lambda_{i,t+1} = \beta^\theta  e^{-\theta \kappa s_{i,t}-\gamma \left(\Delta s_{i,t+1} + \Delta y_{t+1} + \hat{x}_{t+1}\right)}  \left( \frac{n_{i,t+1}+1}{n_{i,t}} \right)^{\theta-1}, \qquad 
    \ell_{i,t+1} = \frac{\sigma_{x,t}}{\sigma_{x,i,t}} e^{-\frac{(\hat{x}_{t+1} - \overline{x}_{i,t})^2}{2\sigma_{x,i,t}^2} + \frac{(\hat{x}_{t+1} - \overline{x}_t)^2}{2\sigma_{x,t}^2}},
\end{equation}\normalsize
$\overline{x}_t = \mu_x + \rho_x (\hat{x}_{t}-\mu_x)$, and $\overline{x}_{i,t} = \mu_{x,i} + \rho_{x,i} (\hat{x}_{t}-\mu_{x,i}) + v_{i,t}$.

\subsection{Extension: Stochastic Volatility}\label{app:sv_extension}

\q{SV}{
    In this subsection, we extend the quantitative model to allow for stochastic volatility, following \citet{bansal2004risks}. We consider two cases. In the first case, investors share common beliefs about the volatility process and volatility follows a standard variance-level stochastic volatility (SV) process. In the second case, investors have heterogeneous beliefs about volatility. In this specification, productivity growth remains iid under the objective measure, so there is no objective stochastic volatility, but investors perceive time-varying volatility.
    
    We also consider two correlation structures: (i) volatility shocks that are orthogonal to productivity and belief shocks, and (ii) volatility shocks that are negatively correlated with belief shocks. The latter captures episodes in which a negative sentiment shock is associated with both lower expected productivity growth and higher (perceived or objective) volatility.
    
    Table~\ref{tab:uncond_moments_sv_rhoqv0_rhoxq_grid} reports the results. The first column reproduces the high labor-elasticity baseline from Table~\ref{tab:uncond_moments}. The next two columns introduce stochastic volatility under common beliefs. Stochastic volatility increases the equity premium and return volatility relative to the baseline, while leaving consumption growth, dividend growth, and hours close to their benchmark values. Introducing a negative correlation between volatility and belief shocks reduces interest-rate volatility, in line with the upper bound of the empirical interval.
    
    The final two columns allow for heterogeneous beliefs about volatility. Relative to the baseline, equity premia and return volatility again increase modestly. Interest-rate volatility is somewhat higher under heterogeneous volatility beliefs, but this effect is attenuated when volatility shocks are negatively correlated with belief shocks.
    
    Overall, stochastic volatility acts primarily as an amplification mechanism for the equity premium and return volatility. Importantly, the main qualitative features of the baseline model remain intact. This extension shows that the core results are not driven by the absence of stochastic volatility and that introducing volatility shocks does not overturn the model's implications for macroeconomic quantities.
}
\begin{table}[t]
    \centering
    \caption{Unconditional moments: Extension with stochastic volatility}
    \label{tab:uncond_moments_sv_rhoqv0_rhoxq_grid}
    \resizebox{\textwidth}{!}{%
    \renewcommand{\arraystretch}{1.75}
    \begin{tabular}{@{}lcccccccccccc@{}}
    \toprule\toprule
    & \multicolumn{2}{c}{\textbf{Baseline}}
    & \multicolumn{4}{c}{\textbf{SV: Common Beliefs}}
    & \multicolumn{4}{c}{\textbf{SV: Heterogeneous Beliefs}}
    & \multicolumn{2}{c}{} \\
    \cmidrule(lr){2-3}\cmidrule(lr){4-7}\cmidrule(lr){8-11}\cmidrule(lr){12-13}
    & \multicolumn{2}{c}{\textbf{High elast.}}
    & \multicolumn{2}{c}{\textbf{\(\rho_{qv}=0.0\)}}
    & \multicolumn{2}{c}{\textbf{\(\rho_{qv}=-0.9\)}}
    & \multicolumn{2}{c}{\textbf{\(\rho_{qv}=0.0\)}}
    & \multicolumn{2}{c}{\textbf{\(\rho_{qv}=-0.9\)}}
    & \multicolumn{2}{c}{\textbf{Data (\(\pm 2\) s.e.)}} \\
    \cmidrule(lr){2-3}\cmidrule(lr){4-5}\cmidrule(lr){6-7}\cmidrule(lr){8-9}\cmidrule(lr){10-11}\cmidrule(lr){12-13}
    \textbf{Variables}
    & \textbf{Mean} & \textbf{Std}
    & \textbf{Mean} & \textbf{Std}
    & \textbf{Mean} & \textbf{Std}
    & \textbf{Mean} & \textbf{Std}
    & \textbf{Mean} & \textbf{Std}
    & \textbf{Mean} & \textbf{Std} \\
    \midrule
    Interest rate           & 1.6 & 6.1 & 1.5 & 6.3 & 1.6 & 6.0 & 1.7 & 6.3 & 1.8 & 6.1 & [0.6, 3.0] & [3.7, 6.0] \\
    Excess returns (equity) & 5.3 & 15.6 & 5.7 & 17.4 & 5.7 & 18.1 & 5.5 & 16.9 & 5.1 & 15.4 & [3.5, 9.3] & [15.6, 21.4] \\
    \midrule
    Consumption growth      & 2.1 & 3.0 & 2.2 & 2.9 & 2.2 & 2.9 & 2.2 & 3.0 & 2.2 & 2.9 & [1.2, 3.1] & [2.0, 3.4] \\
    Dividend growth         & 2.1 & 9.8 & 2.2 & 9.7 & 2.2 & 9.7 & 2.2 & 9.8 & 2.2 & 9.8 & [-0.3, 3.4] & [6.9, 11.2] \\
    \midrule
    Log hours               & 0.6 & 1.9 & 0.5 & 1.9 & 0.6 & 1.9 & 0.5 & 1.9 & 0.6 & 1.9 & -- & [1.7, 2.7] \\
    \bottomrule\bottomrule
    \end{tabular}%
    }
    \vspace{0.35em}
    \parbox{\textwidth}{\scriptsize
    \vspace{0.7em}
    \textit{Note:} Model moments are annualized and reported in percentage points. 
The baseline column corresponds to Model 5 (High elasticity) in Table~\ref{tab:uncond_moments}. 
``SV: Common Beliefs'' uses a Bansal--Yaron style variance-level stochastic volatility specification in which both investors share the same volatility process. 
``SV: Heterogeneous Volatility Beliefs'' assumes no objective stochastic volatility, while one investor perceives volatility to follow a variance-level stochastic volatility process. 
\(\rho_{qv}\) denotes the correlation between volatility shocks and belief shocks. 
Data intervals report the same \(\pm 2\) bootstrap standard-error ranges used in Table~\ref{tab:uncond_moments}.
    }
\end{table}

